% =====================================================================
% CAP. 11 - NIVEL 3 - FONTES DEPENDENTES
% =====================================================================
\blocoaprofundamento

\begin{questao}[3]{}
O capacitor abaixo possui $v_C(0^+)=\SI{8}{\volt}$. A fonte dependente é uma VCVS
que impõe $v_b=\mu v_C$, com $\mu=0{,}5$. Determine a resistência equivalente vista
pelo capacitor, a constante de tempo e $v_C(t)$.

\circuitoqq{%
\begin{circuitikz}[scale=1,transform shape,line width=0.8pt]
\draw (0,0) -- (6.5,0);
\draw (0,3) to[C,l={$C=\SI{25}{\micro\farad}$},v^>={$v_C$}] (0,0);
\draw (0,3) to[R,l={$R=\SI{4}{\kilo\ohm}$}] (3.2,3) coordinate(b);
\draw (b) to[american controlled voltage source,l_={$\mu v_C$}] (3.2,0);
\node[right] at (b) {$v_b$};
\end{circuitikz}}{Circuito RC com fonte de tensão controlada por tensão (VCVS).}

\resposta{$R_{\mathrm{eq}}=\SI{8}{\kilo\ohm}$; $\tau=\SI{0.20}{\second}$;
$v_C(t)=8e^{-5t}\ \si{\volt}$}
\resolucao{Aplique uma tensão-teste $v_x$ aos terminais do capacitor. Como a VCVS
impõe $v_b=\mu v_x$, a corrente de teste é
\[
i_t=\frac{v_x-v_b}{R}=\frac{(1-\mu)v_x}{R}.
\]
Logo
$R_{\mathrm{eq}}=v_x/i_t=R/(1-\mu)=4\,\mathrm{k\Omega}/0{,}5=
\SI{8}{\kilo\ohm}$. Assim
$\tau=R_{\mathrm{eq}}C=\SI{0.20}{\second}$. Como o circuito é natural,
$v_C(\infty)=0$ e $v_C(t)=8e^{-t/0{,}2}=8e^{-5t}\,\si{\volt}$.}
\end{questao}

\begin{questao}[3]{}
O indutor do circuito abaixo conduz inicialmente $\SI{3}{\ampere}$. A fonte
dependente é uma CCVS de valor $r_m i_x$, com $r_m=\SI{9}{\ohm}$, e sua polaridade
é indicada no desenho. Determine a resistência equivalente vista pelo indutor,
$\tau$ e $i_L(t)$.

\circuitoqq{%
\begin{circuitikz}[scale=1,transform shape,line width=0.8pt]
\draw (0,0) -- (6,0);
\draw (0,3) to[L,l={$L=\SI{75}{\milli\henry}$},i>^={$i_L$}] (0,0);
\draw (0,3) to[R,l={$R=\SI{6}{\ohm}$},i>^={$i_x$}] (3.3,3)
      to[american controlled voltage source,l_={$r_m i_x$}] (3.3,0);
\end{circuitikz}}{Resposta natural RL com fonte de tensão controlada por corrente (CCVS).}

\resposta{$R_{\mathrm{eq}}=\SI{15}{\ohm}$; $\tau=\SI{5}{\milli\second}$;
$i_L(t)=3e^{-200t}\ \si{\ampere}$}
\resolucao{Com uma fonte-teste aplicada aos terminais do indutor, a queda total na
rede é a soma da queda no resistor e da tensão da CCVS:
\[
v_t=Ri_x+r_mi_x=(R+r_m)i_x.
\]
Assim $R_{\mathrm{eq}}=R+r_m=6+9=\SI{15}{\ohm}$ e
$\tau=L/R_{\mathrm{eq}}=0{,}075/15=\SI{5}{\milli\second}$. Portanto
$i_L(t)=3e^{-t/\tau}=3e^{-200t}\,\si{\ampere}$.}
\end{questao}

\begin{questao}[3]{}
O RLC em paralelo abaixo não possui fonte independente para $t>0$. A VCCS fornece
uma corrente $g v$, com $g=\SI{50}{\milli\siemens}$. Considere
$v(0^+)=\SI{5}{\volt}$ e $i_L(0^+)=0$. Determine a equação diferencial da resposta
natural, classifique o amortecimento e obtenha $v(t)$.

\circuitoqq{%
\begin{circuitikz}[scale=1,transform shape,line width=0.8pt]
\draw (0,0) -- (8,0);
\draw (0,3) -- (8,3);
\draw (1.6,3) to[R,l={$R=\SI{20}{\ohm}$}] (1.6,0);
\draw (3.8,3) to[C,l={$C=\SI{10}{\milli\farad}$},v^>={$v$}] (3.8,0);
\draw (6,3) to[L,l={$L=\SI{1}{\henry}$},i>^={$i_L$}] (6,0);
\draw (8,3) to[american controlled current source,l_={$g v$}] (8,0);
\end{circuitikz}}{RLC em paralelo com VCCS modificando o amortecimento.}

\resposta{$\ddot v+10\dot v+100v=0$; subamortecida;
$v(t)=e^{-5t}[5\cos(8{,}66t)-2{,}89\sin(8{,}66t)]\ \si{\volt}$}
\resolucao{A KCL fornece
$C\dot v+v/R+gv+i_L=0$. Derivando e usando $\dot i_L=v/L$:
\[
C\ddot v+\left(\frac1R+g\right)\dot v+\frac{1}{L}v=0.
\]
Com os valores dados,
$\ddot v+10\dot v+100v=0$. Logo
$\alpha=5\,\mathrm{s^{-1}}$ e $\omega_0=10\,\mathrm{rad/s}$, resultando em
$\omega_d=\sqrt{100-25}=\SI{8.66}{\radian\per\second}$ e resposta subamortecida.
Da KCL em $0^+$,
$0{,}01\dot v(0)+5/20+0{,}05\cdot5+0=0$, portanto
$\dot v(0)=-\SI{50}{\volt\per\second}$. Assim
$v=e^{-5t}(A\cos8{,}66t+B\sin8{,}66t)$, com $A=5$ e
$-5A+8{,}66B=-50$, resultando em $B\approx-2{,}89$.}
\end{questao}

\begin{questao}[3]{}
No circuito RLC em série abaixo, a CCVS produz uma tensão $r_m i$ que se soma à
queda resistiva, com $r_m=\SI{30}{\ohm}$. Para
$R=\SI{10}{\ohm}$, $L=\SI{0.1}{\henry}$, $C=\SI{100}{\micro\farad}$,
$v_C(0^+)=\SI{10}{\volt}$ e $i(0^+)=0$, determine $v_C(t)$.

\circuitoqq{%
\begin{circuitikz}[scale=1,transform shape,line width=0.8pt]
\draw (0,0) to[R,l={$R=\SI{10}{\ohm}$}] (2.1,0)
      to[american controlled voltage source,l_={$r_m i$}] (4.2,0)
      to[L,l={$L=\SI{0.1}{\henry}$},i>^={$i$}] (6.3,0)
      to[C,l={$C=\SI{100}{\micro\farad}$},v^>={$v_C$}] (8.5,0)
      -- (8.5,-2) -- (0,-2) -- (0,0);
\end{circuitikz}}{RLC em série com CCVS equivalente a amortecimento adicional.}

\resposta{$v_C(t)=e^{-200t}[10\cos(244{,}95t)+8{,}17\sin(244{,}95t)]\ \si{\volt}$}
\resolucao{A CCVS se comporta, para a equação de malha, como uma queda adicional
$r_mi$. Logo a equação natural é
\[
L\ddot v_C+(R+r_m)\dot v_C+\frac1C v_C=0.
\]
Assim
$\alpha=(R+r_m)/(2L)=40/0{,}2=\SI{200}{\per\second}$ e
$\omega_0=1/\sqrt{LC}\approx\SI{316.23}{\radian\per\second}$, de modo que
$\omega_d\approx\SI{244.95}{\radian\per\second}$. A forma é
$v_C=e^{-200t}(A\cos\omega_dt+B\sin\omega_dt)$. Como $A=10$ e
$i(0)=C\dot v_C(0)=0$, temos $-200(10)+244{,}95B=0$, logo
$B\approx\SI{8.17}{\volt}$.}
\end{questao}

\gabarito[Nível 3 --- fontes dependentes]
